\documentclass[11pt,a4paper]{article}
\usepackage[utf8]{inputenc}
\usepackage[T1]{fontenc}
\usepackage[spanish,es-nodecimaldot]{babel}
\usepackage[margin=2.2cm]{geometry}
\usepackage{amsmath,amssymb,booktabs,longtable,array}
\usepackage{xcolor}
\usepackage[hidelinks]{hyperref}
\usepackage{microtype}
\setlength{\parskip}{5pt}
\setlength{\parindent}{0pt}
\setlength{\emergencystretch}{3em}
\newcommand{\sat}{\operatorname{sat}}
\newcommand{\clip}{\operatorname{clip}}
\newcommand{\sgn}{\operatorname{sgn}}
\newcommand{\code}[1]{\texttt{\detokenize{#1}}}
\title{Generación de efectos hápticos en MiYoke\_PD\_V2.3\\\large Fundamento, trazabilidad, ganancias y límites de salida}
\author{Documento técnico del proyecto}
\date{4 de octubre de 2026}
\begin{document}
\maketitle
\begin{abstract}
Se describe el código actual del puente SimConnect SC6 y del firmware Pico 2: sensación prevista, datos de Prepar3D, ecuaciones, antecedentes de software y relación entre ganancias y PWM. El sistema produce comandos normalizados. No mide todavía fuerza o torque y no dispone de una calibración que permita afirmar valores en N o N\,m. Los rangos numéricos son límites de cálculo y no mediciones del esfuerzo de un Cessna R182. Se diferencia una adaptación de TelemFFB de una implementación propia y de un modelo físico verificable.
\end{abstract}
\tableofcontents
\newpage
\section{Qué se siente y qué efectos existen}
La sensación descrita supone un actuador conectado, con polaridad correcta y capacidad mecánica suficiente. En la etapa declarada del proyecto se observan cálculos y PWM en un LCD; la generación de una señal no prueba que esa sensación haya sido reproducida físicamente.
\begin{longtable}{p{3.1cm}p{6.5cm}p{5cm}}
\toprule Efecto & Sensación prevista en el yugo & Ajuste del menú\\\midrule\endhead
Centrado dinámico & Resistencia al apartar el mando del centro; aumenta con presión aerodinámica. Permanece al sostener una posición. & Centrado pitch; Centrado roll\\
Rigidez en reposo & Resistencia residual cuando la presión es pequeña. Es parte del centrado, no una vibración adicional. & Centro reposo\\
Turbulencia / ráfagas & Empujes variables hacia ambos lados y adelante/atrás asociados a cambios del viento. No es un tono periódico fijo. & Turbul. pitch; Turbul. roll\\
Buffet de pérdida & Vibración creciente al aproximarse al ángulo de pérdida. Portadora triangular de 13 Hz. & Buffet perdida\\
Buffet del tren & Vibración con tren retráctil extendido y velocidad suficiente, a 10 Hz. No incluye el golpe de bloqueo del mecanismo. & Buffet tren\\
Buffet de spoilers & Vibración de 15 Hz al extender spoilers a velocidad suficiente. En el addon R182 analizado su salida es cero. & Buffet spoiler\\
Rodaje & Vibración de 8 Hz al desplazarse sobre el suelo, predominante en roll. Se anula en reposo y en vuelo. & Vib. suelo\\
Amortiguación viscosa & Resistencia opuesta a la velocidad del mando; al mover más rápido aumenta hasta su límite. & Damper\\
Fricción suavizada & Resistencia al movimiento que tiende a una meseta. Se anula en reposo; no simula adherencia estática exacta. & Friccion\\\bottomrule
\end{longtable}
Son ocho familias de efectos, contando rigidez residual dentro del centrado. El menú contiene doce ganancias: la total, dos centrados, reposo, dos turbulencias, suelo, tres buffets, damper y fricción. La vista EFECTOS contiene once interruptores, excluyendo la ganancia total. No se implementan actualmente trim físico, flujo de hélice en el centrado, vibración del motor, golpe de aterrizaje separado, movimiento hidráulico del tren ni fuerza en pedales, aunque TelemFFB pueda incluir funciones semejantes.

\section{Arquitectura, magnitudes y límites comunes}
\subsection{Recorrido de los datos}
Prepar3D calcula el estado del avión. \code{SimConnectBridge/main.cpp} solicita veinte variables, calcula presión normalizada, ráfagas y envolventes; envía tramas SC6 de 26 bytes. \code{SimTelemetry.h} valida rango, CRC y secuencia. \code{SimEffects.h} sintetiza y suma los efectos; \code{MotorDriver.h} convierte el comando a PWM. MobiFlight proporciona los ejes al simulador; el firmware consulta esas posiciones con \code{YOKE_USE_SIMULATION=1}. Los encoders son una alternativa futura, no la fuente actual.

El puente limita los envíos normales a intervalos de al menos 20 ms; el firmware calcula nominalmente cada 2 ms. El CSV del puente se escribe por muestra recibida y no equivale a un registro de paquetes enviados. Un watchdog de 300 ms y la desconexión anulan los efectos. La calibración permite pausa o ausencia de enlace; sus salidas de motor son cero. Ausencia de enlace no prueba que el proceso Prepar3D esté cerrado.

\subsection{Notación y conversión a PWM}
Sean $G\in[0,1]$ la ganancia total, $g_i\in[0,1]$ una ganancia de efecto, $e_i\in\{0,1\}$ su interruptor y $\widetilde g_i=e_i g_i$ la ganancia efectiva. En las fórmulas siguientes se supone activado el efecto; si se desactiva se reemplaza su ganancia por cero. $x_R,x_P\in[-1,1]$ son posiciones normalizadas; $v_R,v_P$ son recorridos centro--extremo por segundo.
\begin{align}
\clip(z,a,b)&=\min(b,\max(a,z)), & \sat(z)&=\clip(z,0,1),\\
\mathcal R(z;a,b)&=\sat\left(\frac{z-a}{b-a}\right),&L&=\frac{220}{255}=0.862745\ldots
\end{align}
La suma por eje es $u_a=\clip(\sum_i u_{a,i},-L,L)$. La polaridad configurada es $d_a=\pm1$; actualmente es positiva en ambos ejes. Con \code{PWM_MAX=255}:
\begin{equation}
 PWM_a=\clip\left(\sgn(d_au_a)\left\lfloor255|u_a|+0.5\right\rfloor,-220,220\right).
\end{equation}
PWM firmado positivo/negativo selecciona la dirección del puente H. Su magnitud es un entero; 220 representa aproximadamente 86.3\% del ciclo útil de escala 255. El driver vuelve a limitar la salida a ese techo. \textbf{No es un límite de corriente ni una garantía térmica del motor.}

Para compactar rangos se define
\begin{equation}\mathcal P(z)=\min\{220,\lfloor255z+0.5\rfloor\},\qquad z\geq0.\end{equation}
Un rango aislado de comando $[-z,z]$ corresponde al intervalo contenido en $[-\mathcal P(z),\mathcal P(z)]$. No necesariamente se recorren todos sus enteros en un ensayo. La suma de efectos puede saturar antes que cada efecto aislado; sus máximos no se suman como márgenes independientes.


\subsection{Explicación sencilla: cómo leer los símbolos comunes}
Una fórmula de este documento combina tres clases de datos: \textbf{entrada del simulador}, \textbf{ajuste del usuario} y \textbf{resultado calculado}. No todas las letras son variables de SimConnect.
\begin{description}
\item[$R$, $P$ y $a$.] $R$ indica roll y $P$ pitch. El subíndice $a$ significa ``el eje que se está calculando''. No son datos recibidos.
\item[$G$ y $g_i$.] $G$ es Ganancia total: 60\% significa $G=0.60$. $g_i$ es la ganancia del efecto $i$, por ejemplo 0.120 para Damper. Ambos vienen del menú, no de Prepar3D.
\item[$e_i$ y $\widetilde g_i$.] El interruptor vale uno si el efecto está activado y cero si está desactivado. Multiplicarlo por la ganancia permite apagar el efecto sin borrar el ajuste.
\item[$u$ y PWM.] $u$ es el comando calculado antes de convertirlo a una salida entera. Su signo selecciona sentido y su magnitud determina PWM. No representa newtons. Por ejemplo, $u=0.10$ produce $\lfloor25.5+0.5\rfloor=26$ PWM.
\item[$\clip$, $\sat$ y $\mathcal R$.] Limitar con $\clip$ impide salir de dos extremos; $\sat$ limita entre cero y uno. $\mathcal R(z;a,b)$ es una rampa: vale cero hasta $a$, crece linealmente y llega a uno en $b$. Por ejemplo, $\mathcal R(125;100,150)=0.5$. En esta función $a,b$ son extremos, no ejes.
\item[$|z|$, $\sgn(z)$ y $\min$.] $|z|$ toma la magnitud sin signo; $\sgn$ devuelve la dirección; $\min$ elige el menor valor. $\lfloor z\rfloor$ es la parte entera hacia abajo. Sumando 0.5 antes se redondea la magnitud al entero más cercano.
\item[$\sum_i$, $L$, $d_a$ y $\mathcal P$.] Se suman los comandos de todos los efectos y se limita a $\pm L=\pm220/255$. $d_a$ corrige el sentido según el cableado/configuración. $\mathcal P(z)$ convierte una magnitud normalizada en el máximo PWM entero permitido. Es una abreviatura matemática, no una variable de SimConnect.
\end{description}
El intervalo $[-M,M]$ significa que la salida puede tener ambos signos y no superar la magnitud $M$. No significa que esté oscilando permanentemente entre esos extremos. Las unidades y la condición de vuelo determinan si se llega a ellos.

\subsection{Por qué todavía no se puede indicar una fuerza en newtons}
El campo de comando o ``F'' del LCD es adimensional. La transmisión eléctrica/mecánica requiere, como mínimo,
\begin{equation}
 V_m=RI+L_m\frac{dI}{dt}+K_e\omega_m,\quad \tau_m=K_t I,\quad
 \tau_R\simeq3\eta_RK_tI_R,\quad F_P\simeq\frac{19\eta_PK_tI_P}{r_{\rm ef}}.
\end{equation}
Los reductores declarados son 3:1 y 19:1. Faltan $K_t$, resistencia, eficiencia, radio o paso de transmisión, pérdidas y mediciones. ``Motor 775, 24 V, fuente de 20 A'' no determina estas magnitudes. Por ello el rango físico en N y N\,m es \textbf{no determinado} para todos los efectos. La ecuación de motor es un modelo estándar de motor DC, no una fórmula extraída de TelemFFB \cite{motor}. Las cotas de comando y PWM sí se determinan en este documento.


\subsection{Explicación sencilla: del PWM a la fuerza física}
La tensión aplicada $V_m$ se reparte entre resistencia del bobinado ($RI$), inductancia ($L_m\,dI/dt$) y tensión generada por el motor al girar ($K_e\omega_m$). $I$ es corriente; $R$ resistencia eléctrica, que en esta ecuación no significa roll; $L_m$ inductancia; $\omega_m$ velocidad angular del motor y $K_e$ su constante eléctrica. Ninguna de estas magnitudes se recibe actualmente desde SimConnect.

El torque del motor $\tau_m$ depende de la corriente mediante $K_t$, su constante de torque. Las reducciones 3 y 19 multiplican aproximadamente ese torque y $\eta_R,\eta_P$ representan pérdidas de eficiencia. En pitch, dividir por el radio efectivo $r_{\rm ef}$ convierte torque en fuerza lineal $F_P$. Si la transmisión no es una polea debe usarse la relación mecánica correspondiente. Por eso un PWM conocido no permite deducir por sí solo una fuerza: falta saber cuánta corriente circula y cómo transmite el esfuerzo el mecanismo.

\section{Variables de Prepar3D obtenidas por SimConnect}
Los nombres y unidades siguientes son los solicitados efectivamente en \code{main.cpp}, no todos los datos posibles de TelemFFB. El SDK documenta las variables \cite{sdk}.
\small
\begin{longtable}{p{7cm}p{2.1cm}p{5.5cm}}
\toprule Variable solicitada & Unidad & Utilización\\\midrule\endhead
\code{DYNAMIC PRESSURE} & psf & Centrado; escala de ráfagas. Conversión a Pa.\\
\code{AIRSPEED INDICATED} & knots & Envolventes de pérdida, tren y spoilers.\\
\code{AIRCRAFT WIND X} & knots & Componente lateral para ráfagas.\\
\code{AIRCRAFT WIND Y} & knots & Componente vertical para ráfagas.\\
\code{AIRCRAFT WIND Z} & knots & Componente longitudinal para ráfagas.\\
\code{SIM ON GROUND} & Bool & Activa rodaje y suprime ráfagas/buffets en tierra.\\
\code{IS SLEW ACTIVE} & Bool & Inhibe cálculo activo en modo desplazamiento.\\
\code{INCIDENCE ALPHA} & degrees & Ángulo de ataque para pérdida.\\
\code{STALL ALPHA} & degrees & Ángulo de pérdida de referencia.\\
\code{TRAILING EDGE FLAPS LEFT PERCENT} & fraction & Promedio de flaps para pérdida.\\
\code{TRAILING EDGE FLAPS RIGHT PERCENT} & fraction & Ídem. Se solicita percent over 100.\\
\code{DESIGN SPEED VS0} & knots & Referencia de velocidad con flaps.\\
\code{DESIGN SPEED VS1} & knots & Referencia de velocidad limpia.\\
\code{GEAR LEFT POSITION} & fraction & Extensión usada por buffet del tren.\\
\code{IS GEAR RETRACTABLE} & Bool & Habilita buffet del tren retráctil.\\
\code{SPOILERS LEFT POSITION} & fraction & Promedio de extensión de spoilers.\\
\code{SPOILERS RIGHT POSITION} & fraction & Ídem.\\
\code{GROUND VELOCITY} & knots & Intensidad del rodaje.\\
\code{YOKE X POSITION} & Position & Posición roll para centrado y derivación de velocidad.\\
\code{YOKE Y POSITION} & Position & Posición pitch para centrado y derivación de velocidad.\\\bottomrule
\end{longtable}
\normalsize
El viento está en ejes del avión, no en ejes geográficos. Las fracciones usan 0--1, no 0--100. Además se reciben los eventos de sistema \code{Pause} y \code{Sim}; no son fuerzas ni filas adicionales de la definición de veinte variables. No se solicita velocidad mecánica del yugo, corriente ni fuerza. Los datos inválidos de un buffet opcional anulan ese buffet; los fallos de datos esenciales desarman el envío.

\section{Centrado dinámico y rigidez de reposo}
\textbf{Sensación:} el mando trata de volver a su centro geométrico. La rigidez aumenta con presión; la fuerza persiste si la posición no es cero. No existe desplazamiento del centro por trim en la implementación actual.
\begin{align}
 q&=47.88025898\,q_{\rm psf}, &q_{\rm ref}&=\tfrac12(1.225)(120\cdot0.514444444)^2\simeq2334.24\;\mathrm{Pa},\\
 p&=\clip(q/q_{\rm ref},0,4),& k^*&=r+(1-r)p,\\
 k_n&=k_{n-1}+\frac{\Delta t}{0.100+\Delta t}(k^*-k_{n-1}),&
 u_{C,a}&=-Gg_{C,a}k_nx_a.
\end{align}
$\Delta t$ se expresa en segundos y $r$ es Centro reposo. $120$ kt es una referencia de escala configurable con \code{--reference-knots}; no es VNE. Aumentar presión por encima de la referencia sigue aumentando rigidez hasta $p=4$, límite de representación SC6. Con $r=0.15$, $0.15\leq k\leq3.55$ en régimen estacionario. Al reiniciar se parte de cero y se filtra.

\textbf{Ganancias:} Centrado pitch (default 0.800), Centrado roll (0.400) y Centro reposo (0.150). Para presión fija, $|u_{C,a}|\leq Gg_{C,a}k$ y el máximo PWM aislado es $\mathcal P(Gg_{C,a}k)$. A $p=0$, el rango es $\pm\mathcal P(Gg_{C,a}r)$; a $p=1$, $\pm\mathcal P(Gg_{C,a})$; en todo el rango SC6, $\pm\mathcal P[Gg_{C,a}(4-3r)]$.

Centro reposo no posee un PWM independiente. Su derivada es $\partial k/\partial r=1-p$: subirlo aumenta la rigidez por debajo de la referencia, no la cambia en $p=1$ y la reduce por encima. Desactivarlo equivale a usar $r=0$, manteniendo su valor guardado. Esta consecuencia de la interpolación debe explicarse al calibrar.

\textbf{Procedencia:} principio de resorte proporcional a presión y deflexión, compatible con $q=\rho V^2/2$ \cite{nasa}. La referencia local es \code{MsfsXpFlightControlsMixIn.py}, métodos \code{_calculate_dynamic_pressures} (línea 437), \code{_calculate_control_coefficients} (563) y \code{_calculate_joystick_spring_coefficients} (857). TelemFFB utiliza presión, coeficientes y un resorte de dispositivo; también incluye flujo de hélice, offsets y trim que aquí no se trasladaron. La normalización a 120 kt, el residual, el filtro y SC6 son decisiones de este proyecto, no parámetros identificados del R182. Implementación: \code{main.cpp} y \code{SimEffects.h::update/command}.


\subsection{Explicación sencilla}
El centrado se calcula en cuatro pasos: leer presión, compararla con una referencia, obtener una rigidez y multiplicarla por cuánto se desplazó el mando.
\begin{description}
\item[$q_{\rm psf}$ y $q$.] \textbf{Entrada SimConnect:} \code{DYNAMIC PRESSURE}, solicitada en psf. $q_{\rm psf}$ es el número recibido; $q$ es ese mismo dato convertido a pascales multiplicando por 47.88025898. No se lee $q$ de una variable diferente ni se calcula aquí a partir de IAS.
\item[$q_{\rm ref}$.] Es una presión de referencia calculada por el puente. 1.225 es la densidad de referencia en kg/m$^3$; 120 es la velocidad elegida en kt y 0.514444444 convierte kt a m/s. No son tres mediciones del simulador. La ecuación física $q=\rho V^2/2$ explica la referencia: $\rho$ es densidad del aire y $V$ velocidad respecto al aire; el puente utiliza directamente la presión recibida para el cálculo operativo.
\item[$p$.] Cociente adimensional $q/q_{\rm ref}$, limitado entre 0 y 4. $p=0.5$ significa la mitad de la presión de referencia, no la mitad de la velocidad.
\item[$r$ y $k^*$.] $r$ es la ganancia \textbf{Centro reposo}. $k^*$ es la rigidez deseada: parte de $r$ a presión cero y agrega $(1-r)p$ al aumentar la presión. No son SimVars.
\item[$k_n$, $k_{n-1}$ y $\Delta t$.] $k_n$ es la rigidez suavizada actual y $k_{n-1}$ la anterior. $\Delta t$ es el tiempo entre cálculos del firmware, en segundos. El término $\Delta t/(0.100+\Delta t)$ decide cuánto se avanza hacia la rigidez deseada; 0.100 s es el ajuste de suavizado. Así no cambia de golpe al llegar una nueva muestra.
\item[$x_a$.] \textbf{Entrada SimConnect:} \code{YOKE X POSITION} para $x_R$ y \code{YOKE Y POSITION} para $x_P$. Cero es el centro; magnitud uno es el extremo normalizado del simulador. No se mide todavía el desplazamiento mecánico mediante encoders.
\item[$g_{C,a}$ y $u_{C,a}$.] La ganancia es \textbf{Centrado roll} o \textbf{Centrado pitch}. El producto $Gg_{C,a}k_nx_a$ calcula el comando. El signo menos hace que se oponga al desplazamiento: con posición positiva, comando negativo.
\end{description}
\textbf{Ejemplo:} con total 60\%, pitch 0.800, reposo 0.150, $p=0.5$ y $x_P=0.5$, la rigidez estabilizada es $0.575$. El comando es $-0.60\cdot0.800\cdot0.575\cdot0.5=-0.138$, aproximadamente $-35$ PWM. Si soltamos hasta posición cero, el centrado pasa a cero aunque la rigidez siga siendo alta. Si se mantiene el mando desplazado, el centrado persiste.

\subsection{Fuente bibliográfica del centrado y del reposo}
El fundamento físico de la presión dinámica corresponde a NASA Glenn \cite{nasa}; los nombres y unidades de presión y posición se verifican en el SDK de Prepar3D \cite{sdk}. El antecedente de software es TelemFFB, \code{MsfsXpFlightControlsMixIn.py}, métodos \code{_calculate_dynamic_pressures}, \code{_calculate_control_coefficients} y \code{_calculate_joystick_spring_coefficients} \cite{tfcentrado}. La conversión de unidades, normalización, interpolación con Centro reposo y filtro de 100 ms se verifican en el código de MiYoke \cite{miyoke}. No se atribuye a NASA ni a TelemFFB una calibración del 15\% residual para el R182.

\section{Turbulencia y ráfagas}
\textbf{Sensación:} perturbaciones no periódicas en pitch y roll asociadas a variaciones del viento. Se suprime en tierra, pausa, Slew o datos no frescos.
Para cada componente $w_i$ del viento en kt:
\begin{align}
 a&=e^{-\Delta t/1.0},&\mu_i&\leftarrow\mu_i+(1-a)(w_i-\mu_i),&h_i&=w_i-\mu_i,\\
 z_R&=0.70\frac{h_X}{5}p_T\gamma,&
 z_P&=\frac{h_Y+0.25h_Z}{5}p_T\gamma,&p_T&=\clip(p,0,1),\\
 \mathbf z&\leftarrow\frac{\mathbf z}{\max(1,\|\mathbf z\|)},&
 b&=1-e^{-\Delta t/0.08},&\mathbf s&\leftarrow\mathbf s+b(\mathbf z-\mathbf s).
\end{align}
$\gamma$ corresponde a \code{--gust-gain} (0--10, lanzador: 1); el 0.25 se puede modificar con \code{--gust-longitudinal}. Estos ajustes del puente no son ganancias del menú. El firmware aplica:
\begin{equation}
 D(s)=\sgn(s)\max\left(\frac{|\clip(s,-1,1)|-0.02}{0.98},0\right),\qquad
 u_{T,a}=Gg_{T,a}D(s_a).
\end{equation}
Por eje, $|u_{T,a}|\leq Gg_{T,a}$ y PWM máximo $\mathcal P(Gg_{T,a})$; no implica que ambos ejes alcancen simultáneamente sus máximos por el límite vectorial. Defaults: Turbul. pitch=0.700 y Turbul. roll=0.350.

\textbf{Procedencia:} \code{TurbulenceModulator.py::update}, filtrado pasaaltos, mezclas 1/0.25/0.7, normalización vectorial y suavizado. La recurrencia local equivale a $h_n=a(h_{n-1}+w_n-w_{n-1})$. Diferencias: constante temporal explícita, escala fija 5 kt, multiplicación por presión, zona muerta, y signo lateral positivo local frente al negativo de la referencia. TelemFFB usa viento relativo y su propia sensibilidad; no se deben intercambiar unidades ni signos sin comprobación. La mezcla es una aproximación háptica, no un cálculo de bisagra del R182. Implementación: \code{GustModel.h} y \code{SimTelemetry.h::gustCommand}.


\subsection{Explicación sencilla}
Primero se separa la parte cambiante del viento de su valor habitual; luego se transforma esa variación en empujes de roll y pitch y se suaviza la salida.
\begin{description}
\item[$w_X,w_Y,w_Z$.] \textbf{Entradas SimConnect:} \code{AIRCRAFT WIND X}, \code{AIRCRAFT WIND Y} y \code{AIRCRAFT WIND Z}, en kt y ejes del avión. $i$ sólo permite escribir la misma operación para cualquiera de las tres componentes.
\item[$\mu_i$ y $h_i$.] $\mu_i$ es un promedio que el puente va actualizando. $h_i=w_i-\mu_i$ es lo que se aparta de ese promedio. Un viento constante termina dando $h_i$ cercano a cero; un cambio produce una perturbación transitoria.
\item[$a=e^{-\Delta t/1.0}$.] $a$ es un coeficiente de memoria, no el eje de salida. $e$ es la base de la función exponencial, no el interruptor $e_i$. $\Delta t$ es el tiempo entre muestras del puente y 1.0 s fija qué tan rápido sigue el promedio al viento.
\item[$z_R,z_P$.] Son comandos provisionales. Dividir por 5 kt fija la sensibilidad. 0.70 pondera el viento lateral; en pitch se suma la variación vertical y un 25\% de la longitudinal. Son coeficientes del algoritmo, no fuerzas medidas.
\item[$p_T$ y $\gamma$.] $p_T$ reutiliza la presión derivada de \code{DYNAMIC PRESSURE}, pero limitada a uno para este efecto. $\gamma$ es \code{--gust-gain} del puente. La ganancia del menú se aplica después; son ajustes distintos.
\item[$\mathbf z$ y $\|\mathbf z\|$.] El vector reúne los dos comandos provisionales y su norma es $\sqrt{z_R^2+z_P^2}$. Dividir ambos por el mismo máximo limita la intensidad conjunta sin cambiar su dirección relativa.
\item[$b$, $\mathbf s$ y 0.08.] $b$ es el coeficiente del suavizado de salida; 0.08 s es su constante temporal. $\mathbf s$ guarda la salida filtrada. Estos valores los calcula el puente, no vienen de Prepar3D.
\item[$D(s)$, 0.02 y 0.98.] El firmware ignora magnitudes hasta 0.02 y reescala lo que queda dividiendo por $1-0.02=0.98$. Por ejemplo, $s=0.51$ da $D(s)=0.5$.
\item[$g_{T,a}$ y $u_{T,a}$.] Son la ganancia \textbf{Turbul. roll/pitch} y el comando final del efecto en ese eje. Se multiplica la señal filtrada por esa ganancia y por Ganancia total.
\end{description}
Con $D(s_P)=0.5$, total 60\% y Turbul. pitch=0.700, el comando es $0.21$, aproximadamente 54 PWM. El signo cambia con la perturbación calculada; no se trata de una vibración fija de una frecuencia determinada.

\subsection{Fuente bibliográfica de turbulencia}
TelemFFB, \code{TurbulenceModulator.py::update}, es el antecedente de los filtros y de la mezcla de ejes \cite{tfturb}. El SDK documenta las tres variables de viento y la presión \cite{sdk}. La escala de 5 kt, las constantes temporales explícitas, el signo lateral local y la zona muerta se verifican en \code{GustModel.h} y \code{SimTelemetry.h} de MiYoke \cite{miyoke}. Es una adaptación háptica, no una ley experimental de fuerza del Cessna.

\section{Síntesis común de vibraciones}
Los buffets y el rodaje reciben envolventes $A\in[0,1]$, filtradas con $\beta=\Delta t/(0.100+\Delta t)$. La onda triangular usa fase $\phi\in[0,1000)$:
\begin{equation}
 T(\phi)=\begin{cases}0.004\phi-1,&\phi<500,\\3-0.004\phi,&\phi\geq500,\end{cases}
 \qquad \phi_n=(\phi_{n-1}+f\Delta t_{\rm ms})\bmod1000.
\end{equation}
Su media es cero y su amplitud nominal es uno. El muestreo discreto puede no coincidir exactamente con todos los picos. Cada buffet tiene $u_P=Gg\bar A T_f$ y $u_R=\tfrac12Gg\bar A T_f$. Por ello sus cotas aisladas son $\pm\mathcal P(Gg)$ y $\pm\mathcal P(Gg/2)$. La relación 2:1 es un ajuste de este yugo, no una medición aeronáutica.


\subsection{Explicación sencilla: cómo se construye una vibración}
La \textbf{envolvente} $A$ indica qué tan intensa debería ser la vibración entre cero y uno. La \textbf{portadora} $T$ cambia periódicamente de signo para hacer vibrar el mando. Multiplicarlas produce una vibración pequeña cuando $A$ es pequeño y grande cuando se aproxima a uno.
\begin{description}
\item[$\bar A$, $\beta$ y $\Delta t$.] La barra significa envolvente suavizada. En cada ciclo se calcula $\bar A_n=\bar A_{n-1}+\beta(A-\bar A_{n-1})$. $\beta=\Delta t/(0.100+\Delta t)$ y $\Delta t$ está en segundos. Es el mismo tipo de suavizado que el de rigidez.
\item[$\phi$, $f$ y $T_f$.] $\phi$ es la fase interna de la onda, entre 0 y 999; no es un ángulo del avión. $f$ es su frecuencia en Hz. $T_f$ es la triangular generada con esa frecuencia. Aquí $f$ no significa fracción de flap como en la sección de pérdida.
\item[$\Delta t_{\rm ms}$ y módulo 1000.] Para avanzar la fase se utiliza el tiempo del firmware en milisegundos. El módulo hace que la fase vuelva al principio al completar un ciclo. La primera rama sube de $-1$ a $+1$ y la segunda baja de $+1$ a $-1$.
\item[$u_P,u_R$.] Se multiplica la onda por envolvente, ganancia del efecto y ganancia total. Para los tres buffets, roll recibe la mitad del comando pitch. Esa relación es un ajuste del dispositivo.
\end{description}
Ni la fase ni la onda triangular se reciben por SimConnect: se generan en la Pico. SimConnect aporta las condiciones que determinan la envolvente de cada efecto.

\section{Buffet de pérdida}
\textbf{Sensación:} temblor de 13 Hz que aumenta al acercarse a la pérdida. No es la simulación de un stick shaker instalado en el R182.
\begin{align}
 f&=(f_L+f_R)/2,&\alpha_0&=\alpha_s(0.5+0.2f),\\
 V_s(f)&=V_{s1}+(V_{s0}-V_{s1})f,&
 A_s&=\sat\left[\mathcal R(\alpha;\alpha_0,\alpha_s)\min\left\{(IAS/V_s(f))^2,1.5\right\}\right].
\end{align}
Se exige avión en vuelo y datos válidos: AoA entre $-90$ y $90$ grados, ángulo de pérdida entre 1 y 45, flaps entre 0 y 1 y velocidades de diseño entre 10 y 400 kt. Datos inválidos producen cero. En los registros del addon $\alpha_s=16.1^\circ$, por lo que el inicio es $8.05^\circ$ limpio y $11.27^\circ$ con flap completo.

\textbf{Ganancia y rango:} Buffet perdida, default 0.240. Usar las cotas comunes de buffet, con envolvente menor que uno fuera de la condición máxima.

\textbf{Procedencia:} \code{BuffetingEffectMixIn.py::ac_update_buffeting} (línea 93), rama con StallAoA conocido. La interpolación de flaps, umbral AoA y factor de velocidad proceden de ese enfoque. La copia de TelemFFB usa TAS en el factor de flujo; el puente local usa IAS, por lo que no son numéricamente equivalentes a toda altitud. TelemFFB tiene modos Classic y Dynamic y dirección modulada; el firmware usa una sola triangular de 13 Hz. Además satura la envolvente antes de aplicar la ganancia, mientras la referencia incorpora intensidad antes de la saturación final. La frecuencia y umbral no están validados como espectro real R182. Implementación: \code{BuffetModel.h::calculate} y \code{SimEffects.h}.


\subsection{Explicación sencilla}
La intensidad crece cuando el ángulo de ataque se acerca a la referencia de pérdida. La posición de flaps modifica desde qué ángulo empieza y qué velocidad se usa como referencia.
\begin{description}
\item[$f_L,f_R,f$.] \textbf{SimConnect:} \code{TRAILING EDGE FLAPS LEFT PERCENT} y \code{TRAILING EDGE FLAPS RIGHT PERCENT}, solicitadas como fracciones 0--1. $f$ es su promedio calculado; no son grados de flap.
\item[$\alpha,\alpha_s$.] \textbf{SimConnect:} \code{INCIDENCE ALPHA} y \code{STALL ALPHA}, en grados. Son el ángulo actual y el ángulo de pérdida suministrado por el modelo.
\item[$\alpha_0$.] Ángulo de inicio del buffet calculado por el puente. Es el 50\% de $\alpha_s$ con flaps limpios y el 70\% con flaps completos. No existe una SimVar que entregue directamente este umbral en la implementación.
\item[$V_{s0},V_{s1},IAS$.] \textbf{SimConnect:} \code{DESIGN SPEED VS0}, \code{DESIGN SPEED VS1} y \code{AIRSPEED INDICATED}, todas en kt. Son referencias de configuración y velocidad indicada actual; las velocidades de diseño no son una medición instantánea de la pérdida con el peso actual.
\item[$V_s(f)$.] Referencia interpolada: empieza en $V_{s1}$ y se acerca a $V_{s0}$ al extender flaps. Se calcula localmente.
\item[$\mathcal R(\alpha;\alpha_0,\alpha_s)$.] Vale cero antes del inicio, aumenta linealmente y llega a uno en el ángulo de pérdida. Representa cuánto se ha avanzado dentro de la zona de buffet.
\item[$(IAS/V_s(f))^2$ y 1.5.] Escalan la envolvente con velocidad. La división elimina las unidades; elevar al cuadrado aproxima una dependencia cuadrática. El factor se limita a 1.5 y después la envolvente total se limita a uno.
\item[$A_s$, $\bar A_s$ y $g_s$.] $A_s$ es la envolvente calculada y $\bar A_s$ su versión suavizada. Con $g_s=$ \textbf{Buffet perdida}, se usa $u_{s,P}=Gg_s\bar A_sT_{13}$ y $u_{s,R}=u_{s,P}/2$.
\end{description}
Si el ángulo está a mitad del intervalo entre inicio y pérdida, y $IAS=V_s(f)$, resulta $A_s=0.5$. Con total 60\% y ganancia 0.240, una vez estabilizado el filtro los picos nominales son unos 18 PWM en pitch y 9 en roll. Esto describe el cálculo, no demuestra que ese umbral coincida con el buffet natural del R182.

\subsection{Fuente bibliográfica del buffet de pérdida}
TelemFFB, \code{BuffetingEffectMixIn.py::ac_update_buffeting}, rama con StallAoA conocido, proporciona el antecedente de la envolvente \cite{tfperdida}. Los nombres SimConnect de AoA, flaps y velocidades de diseño se verifican en el SDK \cite{sdk}. La elección de IAS, la saturación previa a la ganancia y la portadora triangular local se verifican en \code{BuffetModel.h} y \code{SimEffects.h} \cite{miyoke}. Estos documentos no acreditan que 13 Hz sea una frecuencia medida en el R182.

\section{Buffet del tren retráctil}
\textbf{Sensación:} vibración de 10 Hz, predominante en pitch, con el tren desplegado durante el vuelo.
\begin{equation} A_g=g_{\rm ext}\mathcal R(IAS;100,150),\qquad u_{g,P}=Gg_g\bar A_g T_{10},\quad u_{g,R}=\tfrac12u_{g,P}.\end{equation}
$g_{\rm ext}$ es la posición del tren izquierdo, no la ganancia. Requiere avión en vuelo, tren retráctil y extensión válida 0--1. Las velocidades son kt; por debajo de 100 no hay envolvente y a partir de 150 llega a su máximo para tren completamente extendido. Esos umbrales no son VLE ni autorización para ensayar fuera de los límites del avión.

\textbf{Ganancia:} Buffet tren, default 0.160; rangos según la síntesis común. \textbf{Referencia:} \code{MotionEffectsMixIn.py::ac_update_landing_gear} (215), rampa de velocidad por extensión, y frecuencia 10 Hz. La copia de TelemFFB usa senoides, umbral adicional de extensión y tratamiento de unidades propio; aquí se usa triangular y extensión izquierda sin ese umbral. El movimiento hidráulico de 150 Hz y el clunk de TelemFFB no se implementaron. Los números 100/150 se interpretan localmente en kt; no se afirma equivalencia numérica con cualquier perfil TelemFFB. Implementación: \code{BuffetModel.h} y \code{SimEffects.h}.


\subsection{Explicación sencilla}
Se multiplica cuánto tren está extendido por cuánto avanzó la velocidad dentro del intervalo elegido para el buffet.
\begin{description}
\item[$g_{\rm ext}$.] \textbf{SimConnect:} \code{GEAR LEFT POSITION}, entre cero y uno. Es extensión física representada por el simulador, \textbf{no una ganancia del menú}, aunque se utilice la letra $g$.
\item[$IAS$ y $\mathcal R(IAS;100,150)$.] IAS es \code{AIRSPEED INDICATED}, en kt. La rampa es cero a 100 kt o menos, 0.5 a 125 kt y uno a 150 kt o más. Los números 100/150 son parámetros locales del efecto.
\item[$A_g,\bar A_g$.] Son la envolvente resultante y su versión suavizada. Por ejemplo, medio tren y media rampa producen $A_g=0.25$.
\item[$g_g,G,T_{10}$.] $g_g$ es \textbf{Buffet tren}; $G$ es Ganancia total; $T_{10}$ es la triangular de 10 Hz generada en la Pico. No son datos del simulador.
\item[$u_{g,P},u_{g,R}$.] Comandos de este buffet: pitch recibe el producto completo y roll la mitad.
\end{description}
También se consulta \code{IS GEAR RETRACTABLE} para confirmar tren retráctil y \code{SIM ON GROUND} para exigir que el avión esté en vuelo. Con tren completo e IAS=125 kt, la envolvente es 0.5. Con total 60\% y Buffet tren=0.160, los picos nominales estabilizados son unos 12 PWM pitch y 6 roll.

\subsection{Fuente bibliográfica del buffet del tren}
TelemFFB, \code{MotionEffectsMixIn.py::ac_update_landing_gear}, es el antecedente de la envolvente por velocidad y extensión \cite{tfmov}. El SDK define \code{GEAR LEFT POSITION}, \code{IS GEAR RETRACTABLE} y la velocidad indicada \cite{sdk}. La interpretación local de 100/150 como kt y la onda triangular se verifican en el proyecto \cite{miyoke}; no deben presentarse como límites aeronáuticos certificados.

\section{Buffet de spoilers}
\textbf{Sensación:} vibración de 15 Hz por la perturbación asociada a spoilers extendidos.
\begin{equation} A_{sp}=\tfrac12(s_L+s_R)\mathcal R(IAS;80,140),\qquad u_{sp,P}=Gg_{sp}\bar A_{sp}T_{15},\quad u_{sp,R}=\tfrac12u_{sp,P}.\end{equation}
Se exige avión en vuelo y ambas fracciones válidas. \textbf{Ganancia:} Buffet spoiler, default 0.160. Rango genérico según las cotas comunes. \textbf{Para el R182 instalado, el rango esperado es cero}: el addon declara límite de spoilers cero y los registros aportados contienen valores inválidos que el puente rechaza.

\textbf{Referencia:} \code{MotionEffectsMixIn.py::ac_update_spoilers} (321): extensión promedio y rampa de velocidad; los umbrales de la referencia se expresan con conversión kt a m/s. TelemFFB combina senoides de 14, 15, 16 y 18 Hz y condiciones adicionales; el proyecto conserva una triangular de 15 Hz. No se implementa el movimiento separado de spoilers. Implementación: \code{BuffetModel.h} y \code{SimEffects.h}.


\subsection{Explicación sencilla}
La intensidad depende de cuánto spoiler está extendido y de la velocidad. No depende del movimiento del yugo.
\begin{description}
\item[$s_L,s_R$.] \textbf{SimConnect:} \code{SPOILERS LEFT POSITION} y \code{SPOILERS RIGHT POSITION}, entre cero y uno. Su promedio describe la extensión conjunta usada por el algoritmo.
\item[$IAS$ y 80/140.] IAS es \code{AIRSPEED INDICATED}, en kt. Los umbrales 80 y 140 son parámetros del algoritmo: a 110 kt la rampa vale 0.5.
\item[$A_{sp},\bar A_{sp}$.] Envolvente calculada y filtrada. Si el promedio de extensión es 0.5 y la rampa también 0.5, $A_{sp}=0.25$.
\item[$g_{sp},G,T_{15}$.] Son \textbf{Buffet spoiler}, Ganancia total y triangular de 15 Hz. La Pico aplica $u_{sp,P}=Gg_{sp}\bar A_{sp}T_{15}$ y entrega la mitad en roll.
\end{description}
\code{SIM ON GROUND} debe indicar vuelo. Si las posiciones son inválidas, no se inventa una extensión: el resultado se anula. Con total 60\%, ganancia 0.160 y envolvente 0.25, serían aproximadamente 6 PWM pitch y 3 roll en un avión equipado. \textbf{En el addon R182 analizado, este ejemplo es sólo demostrativo: no tiene spoilers y la salida esperada es cero.}

\subsection{Fuente bibliográfica del buffet de spoilers}
TelemFFB, \code{MotionEffectsMixIn.py::ac_update_spoilers}, fundamenta el antecedente de extensión promedio y escalado por velocidad \cite{tfmov}. El SDK define las posiciones de spoilers y la IAS \cite{sdk}. La síntesis simplificada a una triangular de 15 Hz y el tratamiento de entradas inválidas pertenecen a \code{BuffetModel.h} y \code{SimEffects.h} \cite{miyoke}. No se atribuye al R182 un efecto de una superficie que no equipa este addon.

\section{Vibración de rodaje}
\textbf{Sensación:} vibración de amplitud creciente con velocidad de suelo, diez veces mayor en roll que en pitch antes del redondeo.
\begin{align}
 A_t&=\begin{cases}\sat[(V_{\rm suelo}-0.5)/39.5],&\text{en tierra y }V_{\rm suelo}>0.5\;\mathrm{kt},\\0,&\text{otro caso},\end{cases}\\
 u_{t,R}&=\clip(Gg_t\bar A_tT_8,-40/255,40/255),&u_{t,P}&=u_{t,R}/10.
\end{align}
Cuando $A_t=0$ se borra inmediatamente la oscilación, sin cola de filtro. Se usa velocidad sobre suelo, no IAS: el viento con el avión detenido no debe generar rodaje.

\textbf{Ganancia:} Vib. suelo, default 0.180. Cotas: $|u_R|\leq\min(Gg_t,40/255)$, $|u_P|\leq\min(Gg_t,40/255)/10$. PWM máximo: $\min(40,\lfloor255Gg_t+0.5\rfloor)$ en roll y $\lfloor\min(255Gg_t,40)/10+0.5\rfloor$ en pitch. No dividir el PWM ya redondeado.

\textbf{Procedencia: implementación propia simplificada.} TelemFFB \code{MotionEffectsMixIn.py::ac_update_runway_rumble} (116) usa compresión de rueda central y diferencia entre ruedas laterales con pasaaltos. El proyecto no replica esa fórmula: usa velocidad, triangular fija, relación 10:1 y techo 40. No se encontró fundamento R182 para esos coeficientes. No atribuirlos a TelemFFB ni presentarlos como rugosidad medida. Implementación: \code{GroundModel.h::envelope} y \code{SimEffects.h}.


\subsection{Explicación sencilla}
El rodaje aumenta su intensidad al aumentar la velocidad sobre el suelo. Usa una onda fija para representar vibración; no está leyendo cada bache de la pista.
\begin{description}
\item[$V_{\rm suelo}$.] \textbf{SimConnect:} \code{GROUND VELOCITY}, en kt. Es velocidad respecto al suelo, diferente de \code{AIRSPEED INDICATED}.
\item[Condición de tierra.] \code{SIM ON GROUND} debe valer uno. Si está en vuelo o la velocidad no supera 0.5 kt, $A_t=0$.
\item[0.5 y 39.5.] 0.5 kt es el umbral para empezar. El divisor 39.5 es $40-0.5$: la rampa llega a uno a 40 kt. Son elecciones del programa.
\item[$A_t,\bar A_t,T_8$.] Son intensidad calculada, intensidad filtrada y onda triangular de 8 Hz. La onda es interna a la Pico, no proviene de la pista simulada.
\item[$g_t,G$ y $40/255$.] $g_t$ es \textbf{Vib. suelo}; $G$ es Ganancia total. El producto se limita a $40/255$ de comando normalizado para que roll no supere 40 PWM.
\item[$u_{t,R},u_{t,P}$ y 10.] El comando pitch es una décima parte del comando roll antes de convertir a enteros. La división por 10 es la relación entre ejes elegida para este efecto, no una reducción mecánica.
\end{description}
Por ejemplo, a 20.25 kt la envolvente vale 0.5. Con total 60\% y Vib. suelo=0.180, una vez estabilizado el filtro, roll alcanza aproximadamente 14 PWM y pitch 1 PWM. Al detenerse se borra inmediatamente el efecto, sin esperar a que desaparezca lentamente el filtro.

\subsection{Fuente bibliográfica del rodaje}
La fórmula aplicada es propia de MiYoke: \code{GroundModel.h::envelope} y \code{SimEffects.h::update/command} \cite{miyoke}. El SDK sustenta el significado de \code{GROUND VELOCITY} y \code{SIM ON GROUND} \cite{sdk}. TelemFFB, \code{MotionEffectsMixIn.py::ac_update_runway_rumble}, se cita como comparación, pues usa compresión de ruedas y no la fórmula local \cite{tfmov}. No hay una fuente experimental externa que valide los valores 8 Hz, 40 kt o 10:1 para este R182; deben identificarse como parámetros empíricos del prototipo.

\section{Amortiguación y fricción}
\subsection{Velocidad utilizada}
No viene directamente de SimConnect: \code{YokeMotion.h::update} deriva la posición recibida, cuantizada por el protocolo a pasos de 0.0001:
\begin{align}
 v_{\rm bruto}&=\clip\left(\frac{x_n-x_{n-1}}{\Delta t},-10,10\right),\\
 v_n&=v_{n-1}+\frac{\Delta t}{0.040+\Delta t}(v_{\rm bruto}-v_{n-1}).
\end{align}
Ante inversión de signo se reinicia la velocidad filtrada antes de actualizar; después de 40 ms con velocidad bruta de magnitud no mayor que 0.02 se fuerza cero. Un hueco de más de 250 ms reinicializa el estimador; un timestamp repetido no crea velocidad nueva. Son decisiones del proyecto. El eje del simulador puede contener zonas muertas y retardos ajenos al mecanismo; no es una medición directa del encoder.


\subsubsection{Explicación sencilla}
Se calcula qué tanto cambió la posición desde la muestra anterior y se divide por el tiempo transcurrido. Es la misma idea de velocidad igual a distancia dividida por tiempo, pero la posición está normalizada.
\begin{description}
\item[$x_n,x_{n-1}$.] Posición actual y anterior obtenidas de \code{YOKE X POSITION} o \code{YOKE Y POSITION}. El subíndice $n$ identifica la muestra, no una unidad de medida.
\item[$\Delta t$.] Tiempo entre muestras nuevas recibidas por el firmware, en segundos. No se deriva una velocidad nueva al reutilizar el mismo paquete.
\item[$v_{\rm bruto}$.] Velocidad calculada antes del filtro. Si la posición pasa de 0 a 0.2 en 0.1 s, resulta 2 recorridos centro--extremo por segundo. No son m/s ni rad/s.
\item[$-10,10$.] Límites numéricos de velocidad; evitan que un salto muy grande genere valores arbitrarios.
\item[$v_n,v_{n-1}$ y 0.040.] Velocidades filtradas actual y anterior. La constante de 0.040 s suaviza los cambios de la derivada. La velocidad que usan los efectos es $v_n$.
\end{description}
El estimador se reinicia ante pausas largas de datos y detecta detención e inversión de sentido. Estos pasos evitan arrastrar una velocidad ficticia. \textbf{No hay una SimVar de velocidad del yugo leída por el puente}: se calcula a partir de las dos posiciones mencionadas.

\subsection{Damper}
\textbf{Sensación:} oposición al movimiento proporcional a su rapidez, incluso pasando por el centro. Definiendo $m=\max(|v|-0.02,0)$, $c_P=1$, $c_R=1/2$:
\begin{equation}u_{d,a}=-\sgn(v_a)Gc_a\min(g_d m_a,0.4).\end{equation}
Esta expresión corresponde al efecto aislado. \textbf{Ganancia:} Damper, default 0.120. Para cualquier ganancia y $|v|\leq10$:
\begin{equation}|u_{d,a}|\leq Gc_a\min(9.98g_d,0.4),\qquad PWM_{\max,a}=\mathcal P[Gc_a\min(9.98g_d,0.4)].\end{equation}
Con $g_d=0.120$ se alcanza el techo cuando $|v|\geq0.02+0.4/0.120\simeq3.353$ recorridos/s. Más ganancia adelanta la saturación; no aumenta indefinidamente la salida máxima.


\subsubsection{Explicación sencilla}
El damper se parece a un amortiguador: ofrece poca resistencia cuando se mueve lento y más cuando se mueve rápido, hasta su techo. La posición por sí sola no lo activa.
\begin{description}
\item[$v_a$ y $m_a$.] $v_a$ es la velocidad filtrada calculada a partir de \code{YOKE X/Y POSITION}. $m_a=\max(|v_a|-0.02,0)$ descuenta una pequeña zona muerta de 0.02 recorridos/s para no responder a movimientos mínimos.
\item[$g_d$.] Es la ganancia \textbf{Damper}. Multiplica esa rapidez útil: con igual velocidad, mayor ganancia significa mayor resistencia antes de saturar.
\item[$\min(g_dm_a,0.4)$.] Limita la contribución de movimiento a 0.4 antes de aplicar la ganancia total y la relación de ejes. Es un límite de software, no una fuerza de 0.4 N.
\item[$G,c_a,-\sgn(v_a)$.] $G$ es Ganancia total. $c_P=1$ y $c_R=0.5$ hacen que roll tenga la mitad de amplitud para igual velocidad. El signo negativo hace que el comando se oponga al movimiento, independientemente de la posición del mando.
\item[9.98 y $PWM_{\max,a}$.] La máxima velocidad permitida es 10; al restar la zona muerta queda $10-0.02=9.98$. Se inserta en la fórmula para hallar una cota de comando y luego $\mathcal P$ la convierte a PWM. No es el PWM que se produce a cualquier velocidad.
\end{description}
Con $v_P=1$, $m_P=0.98$, Damper=0.120 y total 60\%, resulta $u_{d,P}=-0.07056$, aproximadamente $-18$ PWM. Si el movimiento se invierte, cambia el signo; si se detiene, el estimador acaba en cero y desaparece la resistencia.

\subsubsection{Fuente bibliográfica del damper}
La resistencia viscosa proporcional a velocidad se corresponde con el término mecánico $B\omega$ del modelo DC documentado por IIT Kharagpur \cite{motor}; trasladarlo a un mando lineal da $B\dot{x}$. OpenFFBoard, \code{EffectsCalculator.cpp}, es el antecedente para mantener damper separado del resorte \cite{openff}. La zona muerta, normalización, filtro, ganancia y techo exactos son de \code{YokeMotion.h} y \code{SimEffects.h::motionAdjustedForce} \cite{miyoke}. Las posiciones de origen son las SimVars de yugo descritas por el SDK \cite{sdk}. No se copió un coeficiente físico del R182.

\subsection{Fricción suavizada}
\textbf{Sensación:} resistencia que se aproxima a una meseta al crecer la velocidad; cero al detenerse.
\begin{equation}u_{f,a}=-\sgn(v_a)Gc_a\min\left(g_f\frac{m_a}{0.2+m_a},0.4\right).\end{equation}
\textbf{Ganancia:} Friccion, default 0.030. El máximo del factor suavizado es $h=9.98/10.18\simeq0.980354$:
\begin{equation}|u_{f,a}|\leq Gc_a\min(hg_f,0.4),\qquad PWM_{\max,a}=\mathcal P[Gc_a\min(hg_f,0.4)].\end{equation}
No es fricción estática exacta: no retiene una carga con $v=0$. El parámetro 0.2 es un ajuste de regularización en recorridos/s.


\subsubsection{Explicación sencilla}
La fricción suavizada también se opone al movimiento, pero su intensidad se aproxima a un valor casi constante al aumentar la rapidez. Se diferencia del damper porque no crece linealmente durante todo el intervalo previo al límite.
\begin{description}
\item[$v_a$ y $m_a$.] Son los mismos valores calculados a partir de posiciones SimConnect que utiliza el damper. No se recibe una fuerza de fricción desde Prepar3D.
\item[$g_f$.] Es la ganancia \textbf{Friccion}, conservada aunque se desactive el efecto.
\item[$m_a/(0.2+m_a)$.] Es un factor entre cero y casi uno. Vale cero en reposo; con $m_a=0.2$ vale 0.5; con velocidades mayores se aproxima a uno. El 0.2 tiene las mismas unidades de velocidad normalizada que $m_a$ y fija la transición suave.
\item[$0.4,G,c_a,-\sgn(v_a)$.] El 0.4 limita la magnitud del término; $G$ aplica Ganancia total; $c_a$ establece la relación entre ejes y el signo se opone al movimiento. Son parámetros u operaciones del programa.
\item[$h$ y $PWM_{\max,a}$.] Aquí $h=9.98/(0.2+9.98)$ es el mayor factor suavizado permitido por el límite de velocidad. No es el residuo de viento $h_i$ de turbulencia. Se usa sólo para calcular el máximo teórico.
\end{description}
Con $v_P=1$, $m_P=0.98$, Friccion=0.030 y total 60\%, el comando es aproximadamente $-0.01495$, unos $-4$ PWM. Por eso con esas ganancias es normal obtener números pequeños. Al mantener el mando quieto en una posición desplazada, este efecto se anula: no simula la fuerza necesaria para romper una adherencia estática.

\subsubsection{Fuente bibliográfica de la fricción}
OpenFFBoard, \code{EffectsCalculator.cpp}, aporta el antecedente de una condición de fricción independiente del resorte \cite{openff}. La regularización racional $m/(0.2+m)$, la zona muerta y el techo 0.4 son la implementación de MiYoke en \code{SimEffects.h::motionAdjustedForce}, con velocidad de \code{YokeMotion.h} \cite{miyoke}. No se afirma que esa función racional esté copiada de OpenFFBoard o TelemFFB. El SDK identifica las posiciones utilizadas para derivar velocidad \cite{sdk}; no entrega una fuerza de fricción.

\subsection{Suma, fundamento y procedencia}
Con ambos activados el límite de 0.4 se aplica \textbf{a su suma}, no por separado:
\begin{equation}u_{m,a}=-\sgn(v_a)Gc_a\min\left(g_dm_a+g_f\frac{m_a}{0.2+m_a},0.4\right).\end{equation}
Así $u_{m,a}v_a\leq0$: son cargas resistentes en el modelo. La interpretación de amortiguación corresponde a $F_d=-B\dot x$; para fricción se regulariza la ley de Coulomb. No se deducen $B$ o $F_c$ físicos de las ganancias adimensionales.

\textbf{Referencia de implementación:} OpenFFBoard, \code{Firmware/FFBoard/Src/EffectsCalculator.cpp}, casos de spring, damper y friction \cite{openff}. Es antecedente para separarlos; no se copió su ley numérica completa. La función racional, límites, filtrado y factores de eje son implementación local. No atribuir estas fórmulas a un archivo TelemFFB. \code{SimEffects.h::motionAdjustedForce} conserva el nombre histórico, pero actualmente suma cargas independientes y permite oponerse al movimiento con centrado cero. La suma no multiplica ni deforma las portadoras, salvo la saturación final común.


\subsubsection{Explicación sencilla de la suma de resistencias}
Dentro del paréntesis se suman el término de damper $g_dm_a$ y el de fricción $g_fm_a/(0.2+m_a)$. Después se aplica \textbf{una sola vez} el techo 0.4. No se pueden sumar los dos máximos aislados para afirmar que juntos generan 0.8. Finalmente se aplican sentido, ganancia total y factor de eje.

$u_{m,a}v_a\leq0$ expresa que el comando resistente no apunta en el mismo sentido que el movimiento. Es una propiedad del modelo normalizado, no una medición de potencia eléctrica. En la interpretación física $F_d=-B\dot x$, $F_d$ sería una fuerza, $B$ un coeficiente viscoso y $\dot x$ una velocidad física. Esos valores con unidades todavía no están identificados para el mecanismo; la ganancia Damper es adimensional y no debe llamarse directamente $B$ en N\,s/m.

\section{Ganancias y rangos de salida}
\subsection{Nombre exacto, pasos y defaults}
\begin{longtable}{p{4cm}rrp{6cm}}
\toprule Menú & Default & Paso & Alcance\\\midrule\endhead
Ganancia total & 90\% & 5\% & Escala común $G$; máximo 100\%.\\
Centrado pitch & 0.800 & 0.050 & Sólo rigidez pitch.\\
Centrado roll & 0.400 & 0.050 & Sólo rigidez roll.\\
Centro reposo & 0.150 & 0.050 & Interpolación de rigidez de ambos ejes.\\
Turbul. pitch & 0.700 & 0.050 & Sólo ráfagas pitch.\\
Turbul. roll & 0.350 & 0.050 & Sólo ráfagas roll.\\
Vib. suelo & 0.180 & 0.010 & Rodaje roll y un décimo en pitch.\\
Buffet perdida & 0.240 & 0.010 & Pitch y mitad en roll.\\
Buffet tren & 0.160 & 0.010 & Pitch y mitad en roll.\\
Buffet spoiler & 0.160 & 0.010 & Pitch y mitad en roll; cero en este addon.\\
Damper & 0.120 & 0.010 & Resistencia viscosa, dos ejes.\\
Friccion & 0.030 & 0.010 & Resistencia suavizada, dos ejes.\\\bottomrule
\end{longtable}
Todas las ganancias de efecto están entre 0 y 1. Los defaults son los del código; no prueban qué valores quedaron guardados en la placa. Los ensayos aportados usaron total 60\%, por eso se adopta ese valor para los ejemplos. Activar un efecto con ganancia cero sigue dando cero. El interruptor conserva la ganancia y el estado se guarda con CRC. Los límites de compilación, validez y comunicación siguen teniendo prioridad.

\subsection{Ejemplos con defaults y total 60\%}
Los siguientes son máximos aislados redondeados, con estímulo máximo, no salidas constantes. Centrado se evalúa a recorrido completo; buffets con envolvente uno; damper/fricción con velocidad 10. Cada rango es simétrico $[-M,M]$.
\begin{center}\begin{tabular}{lrr}\toprule Efecto/condición & $M_R$ & $M_P$\\\midrule
Centrado, $p=0$, $r=0.15$ & 9 & 18\\
Centrado, $p=1$ & 61 & 122\\
Centrado, $p=4$, $r=0.15$ & 217 & 220\\
Turbulencia, límite de cada eje & 54 & 107\\
Buffet pérdida & 18 & 37\\
Buffet tren & 12 & 24\\
Buffet spoilers, avión equipado & 12 & 24\\
Buffet spoilers, addon R182 & 0 & 0\\
Rodaje & 28 & 3\\
Damper & 31 & 61\\
Fricción & 2 & 5\\\bottomrule\end{tabular}\end{center}
A 100\% total, el límite global permanece 220. El término conjunto damper/fricción tiene máximo 102 PWM pitch y 51 roll; rodaje mantiene 40/4. Los demás efectos dependen de ganancia/estado y del techo global. La tabla no afirma fuerza física proporcional a estos números.

\subsection{Tabla para cada valor de ganancia a total 60\%}
La tabla siguiente recorre $g=0,0.01,\ldots,1$. Para controles con paso 0.05, usar las filas correspondientes al valor seleccionado; las fórmulas anteriores cubren también valores migrados de memoria. Cada celda es el extremo positivo $M$ del intervalo $[-M,M]$ después del redondeo.

L: $\mathcal P(0.6g)$, válido para turbulencia de un eje, pitch de un buffet y centrado a $p=1$, $|x|=1$. B: $\mathcal P(0.3g)$, roll de un buffet. S: rodaje roll/pitch. D: damper roll/pitch a $|v|=10$. F: fricción roll/pitch a $|v|=10$. Envolventes menores, posiciones parciales y velocidades menores reducen la salida. Centro reposo se evalúa mediante su fórmula, no la columna L.
\small
\begin{longtable}{rrrrrr}
\toprule $g$ & L & B & S (R/P) & D (R/P) & F (R/P)\\\midrule\endhead
0.00 & 0 & 0 & 0/0 & 0/0 & 0/0 \\
0.01 & 2 & 1 & 2/0 & 8/15 & 1/1 \\
0.02 & 3 & 2 & 3/0 & 15/31 & 1/3 \\
0.03 & 5 & 2 & 5/0 & 23/46 & 2/4 \\
0.04 & 6 & 3 & 6/1 & 31/61 & 3/6 \\
0.05 & 8 & 4 & 8/1 & 31/61 & 4/7 \\
0.06 & 9 & 5 & 9/1 & 31/61 & 4/9 \\
0.07 & 11 & 5 & 11/1 & 31/61 & 5/10 \\
0.08 & 12 & 6 & 12/1 & 31/61 & 6/12 \\
0.09 & 14 & 7 & 14/1 & 31/61 & 7/13 \\
0.10 & 15 & 8 & 15/2 & 31/61 & 7/15 \\
0.11 & 17 & 8 & 17/2 & 31/61 & 8/16 \\
0.12 & 18 & 9 & 18/2 & 31/61 & 9/18 \\
0.13 & 20 & 10 & 20/2 & 31/61 & 10/19 \\
0.14 & 21 & 11 & 21/2 & 31/61 & 10/21 \\
0.15 & 23 & 11 & 23/2 & 31/61 & 11/22 \\
0.16 & 24 & 12 & 24/2 & 31/61 & 12/24 \\
0.17 & 26 & 13 & 26/3 & 31/61 & 13/25 \\
0.18 & 28 & 14 & 28/3 & 31/61 & 13/27 \\
0.19 & 29 & 15 & 29/3 & 31/61 & 14/28 \\
0.20 & 31 & 15 & 31/3 & 31/61 & 15/30 \\
0.21 & 32 & 16 & 32/3 & 31/61 & 16/31 \\
0.22 & 34 & 17 & 34/3 & 31/61 & 16/33 \\
0.23 & 35 & 18 & 35/4 & 31/61 & 17/34 \\
0.24 & 37 & 18 & 37/4 & 31/61 & 18/36 \\
0.25 & 38 & 19 & 38/4 & 31/61 & 19/37 \\
0.26 & 40 & 20 & 40/4 & 31/61 & 19/39 \\
0.27 & 41 & 21 & 40/4 & 31/61 & 20/40 \\
0.28 & 43 & 21 & 40/4 & 31/61 & 21/42 \\
0.29 & 44 & 22 & 40/4 & 31/61 & 22/43 \\
0.30 & 46 & 23 & 40/4 & 31/61 & 22/45 \\
0.31 & 47 & 24 & 40/4 & 31/61 & 23/46 \\
0.32 & 49 & 24 & 40/4 & 31/61 & 24/48 \\
0.33 & 50 & 25 & 40/4 & 31/61 & 25/49 \\
0.34 & 52 & 26 & 40/4 & 31/61 & 25/51 \\
0.35 & 54 & 27 & 40/4 & 31/61 & 26/52 \\
0.36 & 55 & 28 & 40/4 & 31/61 & 27/54 \\
0.37 & 57 & 28 & 40/4 & 31/61 & 28/55 \\
0.38 & 58 & 29 & 40/4 & 31/61 & 28/57 \\
0.39 & 60 & 30 & 40/4 & 31/61 & 29/58 \\
0.40 & 61 & 31 & 40/4 & 31/61 & 30/60 \\
0.41 & 63 & 31 & 40/4 & 31/61 & 31/61 \\
0.42 & 64 & 32 & 40/4 & 31/61 & 31/61 \\
0.43 & 66 & 33 & 40/4 & 31/61 & 31/61 \\
0.44 & 67 & 34 & 40/4 & 31/61 & 31/61 \\
0.45 & 69 & 34 & 40/4 & 31/61 & 31/61 \\
0.46 & 70 & 35 & 40/4 & 31/61 & 31/61 \\
0.47 & 72 & 36 & 40/4 & 31/61 & 31/61 \\
0.48 & 73 & 37 & 40/4 & 31/61 & 31/61 \\
0.49 & 75 & 37 & 40/4 & 31/61 & 31/61 \\
0.50 & 77 & 38 & 40/4 & 31/61 & 31/61 \\
0.51 & 78 & 39 & 40/4 & 31/61 & 31/61 \\
0.52 & 80 & 40 & 40/4 & 31/61 & 31/61 \\
0.53 & 81 & 41 & 40/4 & 31/61 & 31/61 \\
0.54 & 83 & 41 & 40/4 & 31/61 & 31/61 \\
0.55 & 84 & 42 & 40/4 & 31/61 & 31/61 \\
0.56 & 86 & 43 & 40/4 & 31/61 & 31/61 \\
0.57 & 87 & 44 & 40/4 & 31/61 & 31/61 \\
0.58 & 89 & 44 & 40/4 & 31/61 & 31/61 \\
0.59 & 90 & 45 & 40/4 & 31/61 & 31/61 \\
0.60 & 92 & 46 & 40/4 & 31/61 & 31/61 \\
0.61 & 93 & 47 & 40/4 & 31/61 & 31/61 \\
0.62 & 95 & 47 & 40/4 & 31/61 & 31/61 \\
0.63 & 96 & 48 & 40/4 & 31/61 & 31/61 \\
0.64 & 98 & 49 & 40/4 & 31/61 & 31/61 \\
0.65 & 99 & 50 & 40/4 & 31/61 & 31/61 \\
0.66 & 101 & 50 & 40/4 & 31/61 & 31/61 \\
0.67 & 103 & 51 & 40/4 & 31/61 & 31/61 \\
0.68 & 104 & 52 & 40/4 & 31/61 & 31/61 \\
0.69 & 106 & 53 & 40/4 & 31/61 & 31/61 \\
0.70 & 107 & 54 & 40/4 & 31/61 & 31/61 \\
0.71 & 109 & 54 & 40/4 & 31/61 & 31/61 \\
0.72 & 110 & 55 & 40/4 & 31/61 & 31/61 \\
0.73 & 112 & 56 & 40/4 & 31/61 & 31/61 \\
0.74 & 113 & 57 & 40/4 & 31/61 & 31/61 \\
0.75 & 115 & 57 & 40/4 & 31/61 & 31/61 \\
0.76 & 116 & 58 & 40/4 & 31/61 & 31/61 \\
0.77 & 118 & 59 & 40/4 & 31/61 & 31/61 \\
0.78 & 119 & 60 & 40/4 & 31/61 & 31/61 \\
0.79 & 121 & 60 & 40/4 & 31/61 & 31/61 \\
0.80 & 122 & 61 & 40/4 & 31/61 & 31/61 \\
0.81 & 124 & 62 & 40/4 & 31/61 & 31/61 \\
0.82 & 125 & 63 & 40/4 & 31/61 & 31/61 \\
0.83 & 127 & 63 & 40/4 & 31/61 & 31/61 \\
0.84 & 129 & 64 & 40/4 & 31/61 & 31/61 \\
0.85 & 130 & 65 & 40/4 & 31/61 & 31/61 \\
0.86 & 132 & 66 & 40/4 & 31/61 & 31/61 \\
0.87 & 133 & 67 & 40/4 & 31/61 & 31/61 \\
0.88 & 135 & 67 & 40/4 & 31/61 & 31/61 \\
0.89 & 136 & 68 & 40/4 & 31/61 & 31/61 \\
0.90 & 138 & 69 & 40/4 & 31/61 & 31/61 \\
0.91 & 139 & 70 & 40/4 & 31/61 & 31/61 \\
0.92 & 141 & 70 & 40/4 & 31/61 & 31/61 \\
0.93 & 142 & 71 & 40/4 & 31/61 & 31/61 \\
0.94 & 144 & 72 & 40/4 & 31/61 & 31/61 \\
0.95 & 145 & 73 & 40/4 & 31/61 & 31/61 \\
0.96 & 147 & 73 & 40/4 & 31/61 & 31/61 \\
0.97 & 148 & 74 & 40/4 & 31/61 & 31/61 \\
0.98 & 150 & 75 & 40/4 & 31/61 & 31/61 \\
0.99 & 151 & 76 & 40/4 & 31/61 & 31/61 \\
1.00 & 153 & 77 & 40/4 & 31/61 & 31/61 \\
\bottomrule\end{longtable}\normalsize
Para otra ganancia total, recalcular con $G$ antes de limitar y redondear: no escalar el entero ya saturado. Para un centrado cualquiera usar $\mathcal P(Ggk|x|)$. Para variar Centro reposo a q fija, insertar el valor $r$ en $k=r+(1-r)p$, manteniendo las dos ganancias de centrado. Esto proporciona el rango para cada valor posible sin atribuir un PWM único a una ganancia aislada de su condición de vuelo.

\section{Trazabilidad y alcance de la fundamentación}
La fuente de TelemFFB es la copia local \code{VPforce-TelemFFB-refactor}, dentro del proyecto, contrastada con el repositorio del autor \cite{telem,archivos}. Se leyó el código, no sólo la documentación comercial. No hay metadatos Git que permitan asignar un commit verificado a esa copia. Los números de línea citados corresponden a la copia auditada y pueden cambiar en la rama remota.
\begin{longtable}{p{3cm}p{5.8cm}p{5.7cm}}
\toprule Función local & Antecedente verificable & Tipo de relación\\\midrule\endhead
Centrado & \code{MsfsXpFlightControlsMixIn.py} & Principio adaptado; no traslado completo de trim/propwash.\\
Ráfagas & \code{TurbulenceModulator.py} & Filtrado y mezcla adaptados; escala, signo y presión locales.\\
Pérdida & \code{BuffetingEffectMixIn.py} & Envolvente adaptada; renderizador simplificado.\\
Tren & \code{MotionEffectsMixIn.py} & Rampa por extensión; portadora y condiciones locales.\\
Spoilers & \code{MotionEffectsMixIn.py} & Envolvente adaptada; una triangular en lugar de varias senoides.\\
Rodaje & \code{MotionEffectsMixIn.py}, sólo como comparación & Fórmula actual propia, distinta de compresión de ruedas.\\
Damper / fricción & OpenFFBoard \code{EffectsCalculator.cpp} & Separación conceptual; ecuaciones y filtros locales.\\\bottomrule
\end{longtable}
Los tres MixIn están, respectivamente, en \code{telemffb/sim/msfs_xp/} o \code{telemffb/sim/base/}; TurbulenceModulator está en \code{telemffb/utils/}. Se entrega junto al TEX un manifiesto \code{fuentes_sha256.json} con hashes completos de los archivos auditados, para identificar la evidencia aunque cambie el repositorio. El TEX es autónomo y no necesita ese JSON para compilar.

\subsection{Qué se puede defender en la exposición}
Se puede afirmar que el sistema transforma telemetría en comandos hápticos mediante relaciones físicas y aproximaciones documentadas, con límites y ganancias reproducibles. No se puede afirmar aún que reproduce exactamente las fuerzas de un Cessna R182: faltan coeficientes de bisagra, transmisión real y validación de carga. Citar un algoritmo de TelemFFB acredita un antecedente de software, no una certificación experimental del avión.

Los ensayos CSV aportados permiten reconstruir respuestas esperadas, pero no contienen PWM medido ni ganancias efectivas de la Pico. El LCD de 10 Hz puede ocultar o aliasar portadoras de 8--15 Hz. Para validar salida deben guardarse posición, velocidad, comando, PWM, ganancias y tiempo de la placa; para validar fuerza se necesita además medición de carga/corriente y caracterización del actuador.

\begin{thebibliography}{99}
\bibitem{telem} Walmis / VPforce. \emph{VPforce-TelemFFB}, rama refactor y copia local auditada. \url{https://github.com/walmis/VPforce-TelemFFB}. Archivos y métodos individualizados en cada sección.
\bibitem{archivos} Rutas verificables de los archivos TelemFFB (rama refactor; su contenido remoto puede cambiar):
\begin{itemize}
\item \href{https://github.com/walmis/VPforce-TelemFFB/blob/refactor/telemffb/sim/msfs_xp/MsfsXpFlightControlsMixIn.py}{MsfsXpFlightControlsMixIn.py}: centrado y controles.
\item \href{https://github.com/walmis/VPforce-TelemFFB/blob/refactor/telemffb/utils/TurbulenceModulator.py}{TurbulenceModulator.py}: perturbaciones de viento.
\item \href{https://github.com/walmis/VPforce-TelemFFB/blob/refactor/telemffb/sim/base/BuffetingEffectMixIn.py}{BuffetingEffectMixIn.py}: buffet de pérdida.
\item \href{https://github.com/walmis/VPforce-TelemFFB/blob/refactor/telemffb/sim/base/MotionEffectsMixIn.py}{MotionEffectsMixIn.py}: tren, spoilers y referencia de rodaje.
\end{itemize}
\bibitem{sdk} Lockheed Martin. \emph{Prepar3D v4 SDK: Simulation Variables}. \url{https://www.prepar3d.com/SDKv4/sdk/references/variables/simulation_variables.html}.
\bibitem{nasa} NASA Glenn Research Center. \emph{Dynamic Pressure}. \url{https://www1.grc.nasa.gov/beginners-guide-to-aeronautics/dynamic-pressure/}.
\bibitem{openff} Ultrawipf y colaboradores. \emph{OpenFFBoard}, \code{Firmware/FFBoard/Src/EffectsCalculator.cpp}. \url{https://github.com/Ultrawipf/OpenFFBoard/blob/master/Firmware/FFBoard/Src/EffectsCalculator.cpp}.
\bibitem{motor} IIT Kharagpur, Virtual Labs. \emph{Study of DC Motor, Theory}. \url{https://vlabs.iitkgp.ac.in/ctrl/Exp3/theory.html}.
\bibitem{tfcentrado} Walmis / VPforce. \emph{MsfsXpFlightControlsMixIn.py}, TelemFFB, rama refactor. Métodos de presión dinámica y coeficientes de controles. \href{https://github.com/walmis/VPforce-TelemFFB/blob/refactor/telemffb/sim/msfs_xp/MsfsXpFlightControlsMixIn.py}{Consultar archivo del repositorio}. Copia local identificada por el manifiesto SHA-256.
\bibitem{tfturb} Walmis / VPforce. \emph{TurbulenceModulator.py}, método update. \href{https://github.com/walmis/VPforce-TelemFFB/blob/refactor/telemffb/utils/TurbulenceModulator.py}{Consultar archivo del repositorio}.
\bibitem{tfperdida} Walmis / VPforce. \emph{BuffetingEffectMixIn.py}, método ac\_update\_buffeting. \href{https://github.com/walmis/VPforce-TelemFFB/blob/refactor/telemffb/sim/base/BuffetingEffectMixIn.py}{Consultar archivo del repositorio}.
\bibitem{tfmov} Walmis / VPforce. \emph{MotionEffectsMixIn.py}, métodos ac\_update\_landing\_gear, ac\_update\_spoilers y ac\_update\_runway\_rumble. \href{https://github.com/walmis/VPforce-TelemFFB/blob/refactor/telemffb/sim/base/MotionEffectsMixIn.py}{Consultar archivo del repositorio}.
\bibitem{miyoke} Proyecto MiYoke\_PD\_V2.3. \emph{Código fuente local auditado}: \code{SimConnectBridge/main.cpp}, \code{GustModel.h}, \code{BuffetModel.h}, \code{GroundModel.h}, \code{SimEffects.h}, \code{SimTelemetry.h}, \code{YokeMotion.h}, \code{GainSettings.h}, \code{MotorDriver.h} y \code{mis_defines.h}. Identificación de versiones mediante \code{fuentes_sha256.json}, entregado junto con este documento. Fuente primaria de los parámetros propios, no una validación experimental del avión.

\end{thebibliography}
\end{document}
